\section{Final Validation Comparison and Model Selection}
\label{sec:final_validation_comparison}

The final validation comparison contains one representative of each model class: Ridge and tuned XGBoost using market features, and tuned Temporal GraphSAGE. Table~\ref{tab:finalist-validation-performance} reports their performance before any test result is examined.

\begin{table}[!htbp]
\centering
\footnotesize
\caption{Validation performance of the three finalist models.}
\label{tab:finalist-validation-performance}
\begin{tabular}{lrrrrr}
\toprule
Model & $\mathrm{MAE}_{\mathrm{train}}$ & $\mathrm{MAE}_{\mathrm{val}}$ & $\mathrm{RMSE}$ & $R^2$ & $\overline{\rho}$ \\
\midrule
Temporal GraphSAGE & 0.0813 & 0.0941 & 0.1348 & 0.440 & 0.748 \\
\textbf{XGBoost} & \textbf{0.0836} & \textbf{0.0951} & \textbf{0.1317} & \textbf{0.466} & \textbf{0.748} \\
Ridge & 0.0869 & 0.0959 & 0.1337 & 0.449 & 0.747 \\
\bottomrule
\end{tabular}
\end{table}


\noindent
Temporal GraphSAGE records the lowest validation MAE, but the difference from XGBoost is less than one tenth of a percentage point. XGBoost has the lower RMSE and higher $R^2$, while the two models have essentially identical rank correlations. Ridge remains competitive but is weaker than XGBoost on each reported validation metric, providing modest support for Hypothesis~H3.

Figure~\ref{fig:finalist-validation-deciles} confirms that the three models make errors in the same parts of the outcome distribution. Temporal GraphSAGE has a small average-MAE advantage, but none of the models resolves the sharp increase in error among the most severe realized drawdowns.

\begin{figure}[!htbp]
    \centering
    \includegraphics[width=0.82\textwidth]{04_07_validation_mae_by_realized_drawdown_decile.pdf}
    \caption[Validation errors across the finalist models]{Validation MAE across realized-drawdown deciles for market-only Ridge, market-only XGBoost, and Temporal GraphSAGE.}
    \label{fig:finalist-validation-deciles}
\end{figure}

\noindent
Market-only XGBoost is therefore selected before the test sample is examined. Although Temporal GraphSAGE achieves a marginally lower validation MAE, XGBoost offers comparable ranking performance, stronger RMSE and $R^2$, and substantially lower implementation complexity. The validation evidence therefore provides mixed rather than conclusive support for Hypothesis~H4.
